\section{Atomic unit (\texttt{amo\_unit})}
\textit{Source: \texttt{design/core/amo\_unit/src/}.}

The atomic unit runs the RV32A instructions: \sig{lr.w} (load-reserved),
\sig{sc.w} (store-conditional), and the \sig{amo*.w} read-modify-write
operations (swap, add, and, or, xor, min, max). Each one needs more than one bus
access, so the unit uses a small FSM and stalls the IE stage while it works.

\diagram{amo-fsm}{The AMO FSM. A load-reserved or AMO reads first; an AMO then
writes the modified value; a store-conditional writes only if its reservation is
still valid. \sig{result\_o} returns the old value (LR/AMO) or success/fail
(SC).}{fig:amo-fsm}{0.7}

\subsection{Ports and state}
Inputs from IE: \sig{amo\_op\_i} (which operation), \sig{addr\_i} (the address,
from the ALU), \sig{rs2\_i} (the second operand), and \sig{flush\_i} (a trap is
taken). Bus side: \sig{dbus\_read\_o}/\sig{dbus\_write\_o}/\sig{dbus\_addr\_o}/
\sig{dbus\_writedata\_o} out, \sig{dbus\_readdata\_i}/\sig{dbus\_stall\_i} in.
Outputs to the core: \sig{result\_o} (value for \sig{rd}), \sig{stall\_o} (stall
IE), \sig{active\_o} (the unit is driving the bus), and \sig{pending\_o} (an AMO
is waiting at IE but has not started yet).

State: \sig{state} (\sig{IDLE}/\sig{AMO\_READ}/\sig{AMO\_WRITE}/\sig{DONE}),
\sig{addr\_q}/\sig{rs2\_q}/\sig{op\_q} (latched request), \sig{read\_data\_q} (the
value read), and the reservation \sig{resv\_valid}/\sig{resv\_addr} for LR/SC.

\subsection{How it runs}
On \sig{lr.w} or any \sig{amo*.w} the FSM goes to \sig{AMO\_READ} and reads the
word. For \sig{lr.w} it returns the value and sets the reservation
(\sig{resv\_valid}, \sig{resv\_addr}). For an \sig{amo*.w} it keeps the old value
in \sig{read\_data\_q}, computes ``old OP \sig{rs2}'', and goes to
\sig{AMO\_WRITE} to store it back; \sig{result\_o} is the old value. On \sig{sc.w}
the FSM first checks the reservation: if \sig{resv\_valid} and \sig{resv\_addr ==
addr\_i} it writes and returns 0 (success); otherwise it writes nothing and
returns 1 (fail). Any \sig{sc.w} or \sig{flush\_i} clears the reservation.

\latencynote{The FSM advances only when the bus is free: each step waits on
\sig{\textasciitilde dbus\_stall\_i}. A slow slave just holds the unit in \sig{AMO\_READ} or
\sig{AMO\_WRITE} longer. \sig{pending\_o} matters for traps: if an interrupt
arrives while an AMO is still waiting at IE, the trap saves \sig{pc\_ie} so the
AMO re-runs after \sig{sret} instead of being skipped. The unit also reads the
bus reply through the per-master \sig{amo\_arb\_rvalid} from \sig{dmem\_arb}, so a
PTW reply can never be mistaken for the AMO's own data.}
